\chapter{Cyanosis}

\section*{Definition}
Bluish or purplish discoloration of the skin, nail beds, and mucous membranes caused by increased deoxygenated hemoglobin (5 g/dL or more) in the capillary bed; central cyanosis involves lips and tongue, while peripheral cyanosis is limited to extremities.

\section*{What Happens in Your Body}
Central cyanosis results from reduced arterial oxygen saturation due to cardiopulmonary disease (right-to-left shunting, severe hypoxemia) or abnormal hemoglobin (methemoglobinemia). Peripheral cyanosis arises from narrowing of blood vessels, reduced blood flow, or increased oxygen extraction in tissues despite normal arterial saturation.

\section*{Causes}
\begin{itemize}
  \item Central: Congenital heart disease (tetralogy of Fallot, transposition of great arteries), pulmonary embolus, pneumonia, COPD, ARDS, high-altitude exposure, methemoglobinemia
  \item Peripheral: Cold exposure, Raynaud phenomenon, peripheral artery disease, heart failure, shock, venous obstruction (deep vein thrombosis)
  \item Mixed: Polycythemia vera, obstructive sleep apnea
  \item Toxins: Nitrates, sulfonamides, benzocaine (methemoglobin inducers)
\end{itemize}

\section*{When to See a Doctor}
See your doctor if:
\begin{itemize}
  \item Central cyanosis (lips, tongue, mucous membranes) requires emergency evaluation
  \item Cyanosis accompanied by dyspnea, chest pain, altered mental status, or clubbing
  \item New-onset cyanosis with no clear environmental cause (cold exposure)
  \item Infants with cyanotic spells or failure to thrive
\end{itemize}

\section*{Related Symptoms}
\begin{itemize}
  \item Dyspnea
  \item Clubbing (fingernails)
  \item Chest pain
  \item Edema
  \item Polycythemia
\end{itemize}
\textbf{Important:} Central cyanosis is a medical emergency until proven otherwise; pulse oximetry may overestimate saturation in methemoglobinemia.

\section*{Self-Care / Home Management}
\begin{itemize}
  \item Warm extremities for peripheral cyanosis due to cold exposure; avoid rewarming too rapidly if frostbite is suspected.
  \item Stop smoking and avoid vasoconstrictive drugs (decongestants, caffeine in excess).
  \item Elevate limbs to reduce dependent edema and improve venous return.
  \item For high-altitude cyanosis, descend to lower altitude immediately.
  \item Avoid self-medicating with prescription medications without medical guidance.
\end{itemize}

\section*{How Your Doctor Will Evaluate You}
\begin{itemize}
  \item Differentiate central vs peripheral by examining mucous membranes and assessing response to supplemental oxygen.
  \item Pulse oximetry and arterial blood gas (ABG) with co-oximetry to measure oxygen saturation and exclude methemoglobin.
  \item Chest X-ray, ECG, and echocardiogram for suspected cardiopulmonary cause.
  \item Complete blood count to assess for polycythemia.
\end{itemize}

\section*{Treatment Options}
\begin{itemize}
  \item Supplemental oxygen for central cyanosis with hypoxemia; treat underlying cardiopulmonary condition.
  \item Methylene blue (1-2 mg/kg IV) for symptomatic methemoglobinemia.
  \item Surgical correction for congenital heart disease; anticoagulation for pulmonary embolus.
  \item Vasodilators and calcium channel blockers for Raynaud-associated peripheral cyanosis.
\end{itemize}

\section*{References}
\begin{thebibliography}{99}
\bibitem{cyanosis:ref1} LeMura A, Liles C, et al. ``Cyanosis.'' In: StatPearls. StatPearls Publishing; 2025.
\bibitem{cyanosis:ref2} Bickler PE, Feiner JR, Lipnick MS, et al. ``Effects of acute hypoxia on the pulse oximeter reading.'' Anesthesiology. 2020;133(2):357-367.
\bibitem{cyanosis:ref3} Martin L, Khalil H. ``How much reduced hemoglobin is necessary to generate central cyanosis?'' Chest. 1990;97(1):182-185.
\bibitem{cyanosis:ref4} Coleman E, Simon R. ``Methemoglobinemia.'' Emerg Med Clin North Am. 2025;43(1):77-90.
\bibitem{cyanosis:ref5} Bernstein D. ``Cyanotic congenital heart disease.'' In: Nelson Textbook of Pediatrics. 22nd ed. Elsevier; 2024.
\bibitem{cyanosis:ref6} Marino BS. ``Evaluation and initial management of cyanosis in the newborn.'' UpToDate. Wolters Kluwer; 2025.
\end{thebibliography}
